% GENERATED FILE. Edit canonical lessons, then run npm run generate:books.
% canonical-source-hash: fnv1a64:ce29dd249f50da30
% canonical-lessons: AR-C437-busala, AR-C437-sura, AR-C437-murshid, AR-C437-arshada, AR-C437-khatt-sayr

\chapter{Finding the way}
\label{ch:ar-finding-the-way}
\hlchaptermodality{437}

\begin{chapteropening}
\textbf{By the end of this chapter:} \emph{I can say a compass, speed, a tour guide, to guide and an itinerary.}

Complete the last lesson of chapter 437: I can say a compass, speed, a tour guide, to guide and an itinerary.
\end{chapteropening}

\section[\texorpdfstring{\ar{بوصلة} (būṣala)}{بوصلة (būṣala)}]{\ar{بوصلة} (būṣala) — a compass}

\label{lesson:AR-C437-busala}



\begin{quote}
Before the new one: say the Arabic for a porter, then the Arabic for to travel about.
\end{quote}

\subsection*{You'll want to know: \ar{بوصلة}}
\textbf{\ar{بوصلة}} — \emph{būṣala} — \textquotedblleft{}a compass\textquotedblright{}.

\begin{quote}
\ar{حمل} \ar{أخي} \ar{بوصلة} \ar{في} \ar{رحلته} \ar{إلى} \ar{الصحراء}.
\end{quote}

\begin{quote}
\emph{ḥamala akhī būṣala fī riḥlatihi ilā aṣ-ṣaḥrāʾ.}
\end{quote}

\begin{quote}
My brother took a compass on his trip to the desert.
\end{quote}

\begin{cousinweb}
\textbf{\ar{بوصلة}} (\emph{būṣala}) was borrowed from Italian \emph{bussola} (EtymDB).
\end{cousinweb}

\begin{grammarlens}[title={What you've built}]
One more word for getting around.
\end{grammarlens}

\subsection*{Your turn}
\begin{itemize}
\raggedright
  \item \emph{Say it:} \emph{būṣala}
  \item \emph{Say it:} \emph{būṣala}, once more
  \item \emph{Say it:} \emph{ḥamala akhī būṣala fī riḥlatihi ilā aṣ-ṣaḥrāʾ.} — My brother took a compass on his trip to the desert.
  \item \emph{Recall:} say the Arabic for a porter, then the Arabic for to travel about, then say \emph{būṣala} again
\end{itemize}

\begin{tcolorbox}[breakable,colback=teal!4,colframe=teal!35!black,title={Before you move on}]
What does \ar{بوصلة} mean? (\textquotedblleft{}A compass\textquotedblright{}.) Say it once more.
\end{tcolorbox}
\section[\texorpdfstring{\ar{سرعة} (surʿa)}{سرعة (surʿa)}]{\ar{سرعة} (surʿa) — speed}

\label{lesson:AR-C437-sura}



\begin{quote}
Before the new one: say the Arabic for to travel about, then the Arabic for a compass.
\end{quote}

\subsection*{You'll want to know: \ar{سرعة}}
\textbf{\ar{سرعة}} — \emph{surʿa} — \textquotedblleft{}speed\textquotedblright{}.

\begin{quote}
\ar{كانت} \ar{سرعة} \ar{القطار} \ar{كبيرة}.
\end{quote}

\begin{quote}
\emph{kānat surʿat al-qiṭār kabīra.}
\end{quote}

\begin{quote}
The train's speed was high.
\end{quote}

\emph{bi-surʿa} means quickly.

\begin{cousinweb}
\textbf{\ar{سرعة}} (\emph{surʿa}) is built on the native Arabic root \emph{s-r-ʿ}, \textquotedblleft{}to be quick\textquotedblright{} (Lane).
\end{cousinweb}

\begin{grammarlens}[title={What you've built}]
One more word for getting around.
\end{grammarlens}

\subsection*{Your turn}
\begin{itemize}
\raggedright
  \item \emph{Say it:} \emph{surʿa}
  \item \emph{Say it:} \emph{surʿa}, once more
  \item \emph{Say it:} \emph{kānat surʿat al-qiṭār kabīra.} — The train's speed was high.
  \item \emph{Recall:} say the Arabic for to travel about, then the Arabic for a compass, then say \emph{surʿa} again
\end{itemize}

\begin{tcolorbox}[breakable,colback=teal!4,colframe=teal!35!black,title={Before you move on}]
What does \ar{سرعة} mean? (\textquotedblleft{}Speed\textquotedblright{}.) Say it once more.
\end{tcolorbox}
\section[\texorpdfstring{\ar{مرشد} (murshid)}{مرشد (murshid)}]{\ar{مرشد} (murshid) — a tour guide; a guide}

\label{lesson:AR-C437-murshid}



\begin{quote}
Before the new one: say the Arabic for a compass, then the Arabic for speed.
\end{quote}

\subsection*{You'll want to know: \ar{مرشد}}
\textbf{\ar{مرشد}} — \emph{murshid} — \textquotedblleft{}a tour guide; a guide\textquotedblright{}.

\begin{quote}
\ar{شرح} \ar{المرشد} \ar{تاريخ} \ar{المدينة}.
\end{quote}

\begin{quote}
\emph{sharaḥa al-murshid tārīkh al-madīna.}
\end{quote}

\begin{quote}
The guide explained the city's history.
\end{quote}

\begin{cousinweb}
\textbf{\ar{مرشد}} (\emph{murshid}) is built on the native Arabic root \emph{r-sh-d}, \textquotedblleft{}to follow the right way\textquotedblright{} (Lane). It is the Form IV participle, \textquotedblleft{}one who leads others the right way\textquotedblright{}.
\end{cousinweb}

\begin{grammarlens}[title={What you've built}]
One more word for getting around.
\end{grammarlens}

\subsection*{Your turn}
\begin{itemize}
\raggedright
  \item \emph{Say it:} \emph{murshid}
  \item \emph{Say it:} \emph{murshid}, once more
  \item \emph{Say it:} \emph{sharaḥa al-murshid tārīkh al-madīna.} — The guide explained the city's history.
  \item \emph{Recall:} say the Arabic for a compass, then the Arabic for speed, then say \emph{murshid} again
\end{itemize}

\begin{tcolorbox}[breakable,colback=teal!4,colframe=teal!35!black,title={Before you move on}]
What does \ar{مرشد} mean? (\textquotedblleft{}A tour guide; a guide\textquotedblright{}.) Say it once more.
\end{tcolorbox}
\section[\texorpdfstring{\ar{أرشد} (arshada)}{أرشد (arshada)}]{\ar{أرشد} (arshada) — to guide}

\label{lesson:AR-C437-arshada}



\begin{quote}
Before the new one: say the Arabic for speed, then the Arabic for a tour guide.
\end{quote}

\subsection*{You'll want to know: \ar{أرشد}}
\textbf{\ar{أرشد}} — \emph{arshada} — \textquotedblleft{}to guide\textquotedblright{}.

\begin{quote}
\ar{أرشدنا} \ar{رجل} \ar{إلى} \ar{الفندق}.
\end{quote}

\begin{quote}
\emph{arshadanā rajul ilā al-funduq.}
\end{quote}

\begin{quote}
A man guided us to the hotel.
\end{quote}

Present: \ar{يرشد}, \emph{yurshid}.

It takes \emph{ilā} before the place.

\begin{cousinweb}
\textbf{\ar{أرشد}} (\emph{arshada}) is built on the native Arabic root \emph{r-sh-d}, \textquotedblleft{}to follow the right way\textquotedblright{} (Lane). It is Form IV, \textquotedblleft{}to lead someone the right way\textquotedblright{}.
\end{cousinweb}

\begin{grammarlens}[title={What you've built}]
One more word for getting around.
\end{grammarlens}

\subsection*{Your turn}
\begin{itemize}
\raggedright
  \item \emph{Say it:} \emph{arshada}
  \item \emph{Say it:} \emph{arshada}, once more
  \item \emph{Say it:} \emph{arshadanā rajul ilā al-funduq.} — A man guided us to the hotel.
  \item \emph{Recall:} say the Arabic for speed, then the Arabic for a tour guide, then say \emph{arshada} again
\end{itemize}

\begin{tcolorbox}[breakable,colback=teal!4,colframe=teal!35!black,title={Before you move on}]
What does \ar{أرشد} mean? (\textquotedblleft{}To guide\textquotedblright{}.) Say it once more.
\end{tcolorbox}
\section[\texorpdfstring{\ar{خط} \ar{سير} (khaṭṭ sayr)}{خط سير (khaṭṭ sayr)}]{\ar{خط} \ar{سير} (khaṭṭ sayr) — an itinerary}

\label{lesson:AR-C437-khatt-sayr}



\begin{quote}
Before the new one: say the Arabic for a tour guide, then the Arabic for to guide.
\end{quote}

\subsection*{You'll want to know: \ar{خط} \ar{سير}}
\textbf{\ar{خط} \ar{سير}} — \emph{khaṭṭ sayr} — \textquotedblleft{}an itinerary\textquotedblright{}.

\begin{quote}
\ar{كتب} \ar{المرشد} \ar{خط} \ar{سير} \ar{الرحلة}.
\end{quote}

\begin{quote}
\emph{kataba al-murshid khaṭṭ sayr ar-riḥla.}
\end{quote}

\begin{quote}
The guide wrote out the itinerary of the trip.
\end{quote}

\begin{cousinweb}
\textbf{\ar{خط} \ar{سير}} (\emph{khaṭṭ sayr}) is built from two native Arabic words: \emph{khaṭṭ}, \textquotedblleft{}a line\textquotedblright{}, on the root \emph{kh-ṭ-ṭ} (Lane), and \emph{sayr}, \textquotedblleft{}going\textquotedblright{}, on the root \emph{s-y-r} (Lane): the line a journey follows.
\end{cousinweb}

\begin{grammarlens}[title={What you've built}]
One more word for getting around.
\end{grammarlens}

\subsection*{Your turn}
\begin{itemize}
\raggedright
  \item \emph{Say it:} \emph{khaṭṭ sayr}
  \item \emph{Say it:} \emph{khaṭṭ sayr}, once more
  \item \emph{Say it:} \emph{kataba al-murshid khaṭṭ sayr ar-riḥla.} — The guide wrote out the itinerary of the trip.
  \item \emph{Recall:} say the Arabic for a tour guide, then the Arabic for to guide, then say \emph{khaṭṭ sayr} again
\end{itemize}

\begin{tcolorbox}[breakable,colback=teal!4,colframe=teal!35!black,title={Before you move on}]
What does \ar{خط} \ar{سير} mean? (\textquotedblleft{}An itinerary\textquotedblright{}.) Say it once more.
\end{tcolorbox}
